\subsection{Long connectors for every terminal pair \((2,\allowbreak b)\)}\label{proof:Q5-C086}
Fix the original finite probability space and a real antisymmetric kernel \(|K|\le\allowbreak 1\).
The terminal has residual arms of lengths \(2,\allowbreak b\), \(b\ge\allowbreak 3\).
The internal vertex of the two-edge arm points toward both the terminal and its leaf; the \(b\)-edge arm points entirely toward its leaf.
All \(L\) connector edges point from the terminal toward the receiving root, whose genuine original unit arm points outward.
The undecorated message is
\[
F_b=h_2f_b,\qquad X_{b,L}=T^LF_b,\qquad H_{b,L}=pX_{b,L}.
\tag{Q86.0.1}\label{eq:Q86:0.1}
\]
Let \(E_t,\allowbreak E_r\in[0,\allowbreak 1]\) be the actual equal-pair factors at the terminal and root.
Each pair uses one fully outward \(f_\ell\) arm and one \(h_\ell\) arm with first edge inward and remaining edges outward, or the full reversal of that pair, fixed by the application.
The full message is \(H_D=\allowbreak E_rpT^L(E_tF_b)\).
All directions are fixed from the tree, its lengths and pairings before the host is given.

For every finite admissible model,
\[
\boxed{E H_D^8\le1-d/20000}
\tag{Q86.0.2}\label{eq:Q86:0.2}
\]
for \(b=\allowbreak 3,\allowbreak L\ge\allowbreak 3\), and for every \(b\ge\allowbreak 4,\allowbreak L\ge\allowbreak 4\).
We also prove the two local bounds
\[
\boxed{E[(p \bar f_n)^8]\le e^{253d/32}\quad(n\ge3),}
\tag{Q86.0.3}\label{eq:Q86:0.3}
\]
\[
\boxed{E[h_2f_b]\le(1-d)/2\quad(b\ge5,\ 1/5\le d\le1/3).}
\tag{Q86.0.4}\label{eq:Q86:0.4}
\]
The first is an upper cost for actual bare paths. The second supplies
\(TF_b\le\allowbreak 1-d\) in its specified high domain.

\paragraph{Q86.1 Original finite model}
Delete zero-mass atoms, and write
\[
Jv=(Ev)1,\quad A=J+K,\quad T=J-K=A^*,\quad K^*=-K,
\quad |K(x,y)|\le1,
\]
\[
g=K1,\ p=1+g,\ q=1-g,\ r=Kg,\ s=Kr,\ t=Ks,
\quad f_n=A^n1,\ \bar f_n=T^n1,\ h_n=T f_{n-1}.
\]
Put \(\delta_j=\allowbreak \|K^j\one\|_2^2,\allowbreak d=\allowbreak \delta_1,\allowbreak \theta=\allowbreak 13/32,\allowbreak \kappa=\allowbreak 19/32,\allowbreak \tau=\allowbreak d-\delta_2\).
All expectations, inner products, adjoints, spectral quantities and norms use \(L^2(\mu)\).
The shared finite calculus in Section~\ref{proof:Q5-C081} gives
\[
0\le d\le1/3,\quad \|K\|^2\le\theta,\quad
d^2\le \delta_2\le\theta d,\quad \delta_{j+1}\le\theta \delta_j,
\quad |Kv(x)|\le\|v\|_2.
\tag{Q86.1.1}\label{eq:Q86:1.1}
\]
Cauchy--Schwarz in the spectral measure of \(-K^2\) also gives
\(d\delta_3\ge\allowbreak \delta_2^2\).
The operators \(A,\allowbreak T\) are positive,
\(0\le\allowbreak  f_n,\allowbreak \bar f_n,\allowbreak h_n\le\allowbreak 2\), and \(\E f_n=\allowbreak \E\bar f_n\le\allowbreak 1\).
Actual equal pairs and their full reversals satisfy \(f_\ell h_\ell\le\allowbreak 1\).
Since \(K(x,\allowbreak x)=\allowbreak 0\), \(p(x),\allowbreak q(x)\ge\allowbreak \mu(x)>\allowbreak 0\).
Every finite rooted-tree message is positive by assigning all vertices the root label.

Let \(\beta_1=\allowbreak \beta_2=\allowbreak 0\), and
\[
\beta_j=\sum_{i=1}^{\lfloor(j-1)/2\rfloor}(-1)^{i-1}\delta_i
\quad(j\ge3),\qquad \tau\le\beta_j\le d.
\]
Finite Duhamel expansion gives
\[
S_n:=\sum_{j=0}^nK^j1
=f_n+\sum_{j=3}^n\beta_jf_{n-j},\qquad
S'_n:=\sum_{j=0}^n(-K)^j1
=\bar f_n+\sum_{j=3}^n\beta_j\bar f_{n-j}.
\tag{Q86.1.2}\label{eq:Q86:1.2}
\]
Every retained or discarded path loss is nonnegative.
For \(G_n(z)=\allowbreak 1+z+\cdots+z^{n-1}\), \(G_0=\allowbreak 0\),
\[
\|G_n(K)\|\le1\quad(n=0,1\text{ or }n\ge3),
\qquad\|G_2(K)\|^2\le1+\theta.
\tag{Q86.1.3}\label{eq:Q86:1.3}
\]
The signed exponential estimate from Section~\ref{proof:Q5-C081} states that
\(|W|\le\allowbreak  U,\allowbreak \E W^2\le\allowbreak \theta U^2\) imply, for every \(s\ge\allowbreak 0\),
\[
Ee^{sW}\le\exp\{sEW+(17/64)s^2U^2\}.
\]
In particular, for every real original function \(V\),
\[
\boxed{Ee^{8KV}\le\exp\{8E(KV)+17\|V\|_2^2\}.}
\tag{Q86.1.4}\label{eq:Q86:1.4}
\]
At \(d=\allowbreak 0\), \(g=\allowbreak 0,\allowbreak A\one=\allowbreak T\one=\allowbreak \one\), and leaf integration gives every actual message one.
All divisions below use \(d>\allowbreak 0\).

\paragraph{Q86.2 Square envelopes and their low-domain data}
For \(b=\allowbreak 3\), \(h_2+f_3=\allowbreak 2-d+s\), so
\[
F_3\le(m+R/2)^2,\qquad m=1-d/2,\quad R=s
\quad(0<d\le1/3).
\tag{Q86.2.1}\label{eq:Q86:2.1}
\]
For \(b=\allowbreak 4\), \(h_2+f_4=\allowbreak 2-2d+s+t-dg\), so
\[
F_4\le(m+R/2)^2,\qquad m=1-d,\quad R=s+t-dg
\quad(0<d\le1/4).
\tag{Q86.2.2}\label{eq:Q86:2.2}
\]
Before squaring, \(2m+R\) is a sum of nonnegative original messages.

For \(b\ge\allowbreak 5\), work on \(0<\allowbreak d\le\allowbreak 1/5\), with \(m=\allowbreak 1-9d/8,\allowbreak k=\allowbreak 1/m\).
For \(b=\allowbreak 5\), set \(C_1=\allowbreak d,\allowbreak C_2=\allowbreak d/4+\delta_2,\allowbreak C_3=\allowbreak 0\).
Then \(C_2\le\allowbreak 21d/32\le\allowbreak \beta_3\) and
\(\beta_5+C_1+C_2=\allowbreak 9d/4\).
For \(b\ge\allowbreak 6\), set
\[
\lambda=\frac{9d/4-\beta_b-\beta_{b-1}}
{\beta_{b-2}+(1-d)\beta_{b-3}},\quad
C_1=\beta_{b-1},\quad C_2=\lambda\beta_{b-2},\quad
C_3=\lambda\beta_{b-3}.
\]
The denominator is positive, the numerator is nonnegative, and the full four-term constant is at least
\((4-d)\tau\ge\allowbreak 361d/160>\allowbreak 9d/4\).
Thus \(0\le\allowbreak \lambda\le\allowbreak 1\), each \(C_i\) is no larger than its actual \(\beta\) coefficient,
\(0\le\allowbreak  C_i\le\allowbreak  d\), and \(C_1\ge\allowbreak \tau\).
Set \(C=\allowbreak C_1+C_2+C_3,\allowbreak D=\allowbreak C_2+C_3,\allowbreak E_0=\allowbreak C_3\), and
\[
R=\sum_{j=3}^bK^j1-Cg-Dr-E_0s.
\tag{Q86.2.3}\label{eq:Q86:2.3}
\]
Retain \(\beta_b,\allowbreak C_1f_1,\allowbreak C_2f_2,\allowbreak E_0f_3\) in \eqref{eq:Q86:1.2}.
Using \(f_1=\allowbreak 1+g,\allowbreak f_2=\allowbreak 1+g+r,\allowbreak f_3=\allowbreak 1-d+g+r+s\), their constant loss is
\(\beta_b+C_1+C_2+(1-d)E_0=\allowbreak 9d/4\). Hence
\[
0\le h_2+f_b\le2m+R,\qquad F_b\le F_0:=(m+R/2)^2.
\tag{Q86.2.4}\label{eq:Q86:2.4}
\]
For \(b=\allowbreak 5\), \(E_0=\allowbreak 0\), so no positive amount is retained from \(\beta_2=\allowbreak 0\).

On the original spectral interval \(\xi\in[0,\allowbreak \theta]\), the multiplier of \(R\) relative to \(g\) is
\(-A(\xi)-i\sqrt\xi B(\xi)\), with \(A,\allowbreak B\ge\allowbreak 0\).
For \(b=\allowbreak 5\),
\[
A=d+C_2+\xi-\xi^2,\qquad B=C_2+\xi.
\]
For \(b\ge\allowbreak 6\), let \(N=\allowbreak b-2\ge\allowbreak 4\) and
\(G_N(i\sqrt\xi)=\allowbreak a_f+i\sqrt\xi b_f\). Then
\[
A=C+\xi(a_f-E_0),\quad B=D+\xi b_f,
\quad 1-\xi\le a_f,b_f\le1-\xi+\xi^2,
\quad a_f\ge(1-\xi)b_f.
\tag{Q86.2.5}\label{eq:Q86:2.5}
\]
These alternating finite-sum bounds include \(N=\allowbreak 4,\allowbreak 5\).
For \(N=\allowbreak 2j\),
\(a_f=\allowbreak b_f=\allowbreak \sum_{i=\allowbreak 0}^{j-1}(-\xi)^i\).
For \(N=\allowbreak 2j+1\),
\(a_f-(1-\xi)b_f=\allowbreak [\xi+(-\xi)^j]/(1+\xi)\ge\allowbreak 0\), \(j\ge\allowbreak 2\).
It follows that
\[
A-(1-\xi)B\ge C_1+C_2\xi,\qquad C_1\ge\tau.
\tag{Q86.2.6}\label{eq:Q86:2.6}
\]
The uniform frequency bounds are
\[
\boxed{\sqrt{A^2+\xi B^2}\le u_0:=3/8+9d/5,
\qquad B\le z_0:=\theta+13d/16.}
\tag{Q86.2.7}\label{eq:Q86:2.7}
\]
For \(b=\allowbreak 5\), with the global \(\delta_2\) fixed, \(A,\allowbreak B\) increase nonnegatively with \(\xi\).
First use \(\xi=\allowbreak \theta\), then \(\delta_2\le\allowbreak \theta d\).
The endpoint vector is bounded by
\((\theta-\theta^2,\sqrt\theta\,\theta)
+d(53/32,\sqrt\theta\,21/32)\).
The squared norms of these two vectors are
\(131313/1048576<\allowbreak 9/64\) and \(95621/32768<\allowbreak (9/5)^2\).
Also \(B\le\allowbreak \theta+21d/32\le\allowbreak  z_0\).

For \(b\ge\allowbreak 6\), the exact retained constant yields
\[
C=9d/4-\beta_b+dE_0,\qquad
D=9d/4-\beta_b-\beta_{b-1}+dE_0.
\]
Let \(u(\xi)=\allowbreak \xi-\xi^2+\xi^3\). The frequency bounds are
\[
A\le\overline A(\xi)=u(\xi)+5d/4+\delta_2+E_0(d-\xi),
\quad
B\le\overline B(\xi)=u(\xi)+d/4+2\delta_2+dE_0.
\tag{Q86.2.8}\label{eq:Q86:2.8}
\]
Throughout the frequency interval, \(d,\allowbreak \delta_2,\allowbreak E_0\) are fixed.
Since \(u'\ge\allowbreak 2/3\) and \(E_0\le\allowbreak  d\le\allowbreak 1/5\), both nonnegative upper bounds increase, as does
\(\overline A^2+\xi\overline B^2\).
Use \(\xi=\allowbreak \theta\) before \(\delta_2\le\allowbreak \theta d\), obtaining
\[
A_* =u+(5/4+\theta)d+E_0(d-\theta),\quad
B_* =u+(1/4+2\theta)d+dE_0,\quad
u=10101/32768<5/16.
\]
The expression \(A_*^2+\theta B_*^2\) is convex in \(E_0\in[0,\allowbreak d]\).
At \(E_0=\allowbreak 0\), the constant-vector norm is less than \(3/8\), and the squared \(d\)-vector norm is
\(26229/8192<\allowbreak (9/5)^2\).
At \(E_0=\allowbreak d\), the squared \(d\)-vector norm is
\(16557/8192<\allowbreak (10/7)^2\), and the \(d^2\)-vector norm is below \(6/5\).
The latter endpoint norm is therefore at most
\(3/8+(10/7)d+(6/5)d^2\le\allowbreak 3/8+(9/5)d\), since
\(10/7+6/25<\allowbreak 9/5\).
This bounds every \(E_0\).
Also \(B\le\allowbreak 5/16+17d/16+d^2\le\allowbreak  z_0\), because the difference
\(-3/32+d/4+d^2\) increases on \([0,\allowbreak 1/5]\) and has right endpoint \(-3/800\).

Write \(B_0=\allowbreak -K,\allowbreak g'=\allowbreak -g,\allowbreak k=\allowbreak 1/m,\allowbreak V=\allowbreak \E R^2,\allowbreak D_j=\allowbreak \E(B_0^jR)\), and
\[
a_1=EKR=-E(gR),\quad
\delta=kD_0+k^2V/2,\quad a=ka_1,\quad
\varepsilon=\delta-a,\quad\nu=1+\varepsilon.
\tag{Q86.2.9}\label{eq:Q86:2.9}
\]
The operator \(B_0\) is distinct from the positive scalar \(B\) used below.
For \(b\ge\allowbreak 4\), use the spectral measure of \(g\), with mass \(d\):
\[
V=\int(A^2+\xi B^2),\quad D_0=\int B,\quad a_1=\int A,
\quad D_2=-\int\xi B,
\]
\[
D_{4j+1}=-\int\xi^{2j}A,\quad
D_{4j+2}=-\int\xi^{2j+1}B,\quad
D_{4j+3}=\int\xi^{2j+1}A,\quad
D_{4j+4}=\int\xi^{2j+2}B.
\tag{Q86.2.10}\label{eq:Q86:2.10}
\]
For \(b\ge\allowbreak 5\) on its low domain, \eqref{eq:Q86:2.6} gives
\(a_1\ge\allowbreak \kappa D_0+\tau d\).
Also \(a_1-D_0\le\allowbreak  d^2+\theta^3d\).
For \(b=\allowbreak 5\), this difference is \(d^2-\delta_3\).
For \(b\ge\allowbreak 6\), use
\(A-B=\allowbreak C_1-E_0\xi+\xi(a_f-b_f)\), whose first possible positive end term has degree at least three.
Consequently,
\[
|\varepsilon|\le3d/5,\quad\delta\le6d/5,\quad ku_0\le1
\qquad(b\ge5,\ d\le1/5).
\tag{Q86.2.11}\label{eq:Q86:2.11}
\]
\(\delta/d\le\allowbreak  z_0/m+u_0^2/(2m^2)\).
After multiplication by \(2m^2\), the required positive difference is
\(463/320-(955/128)d+(2601/1600)d^2\), decreasing on the interval with right endpoint \(1577/80000\).
For the upper \(\varepsilon\) bound, use \(a_1\ge\allowbreak \kappa D_0+\tau d\);
the positive difference, multiplied by \(2m^2\), is
\(1867/2560-(64539/20480)d-(118503/51200)d^2\),
with right endpoint \(16519/2560000\).
For its lower bound,
\(-\varepsilon/d\le\allowbreak  k(d+\theta^3)<\allowbreak 3/5\).
Finally \(u_0/m\) increases, and \(ku_0\le\allowbreak 1\) at the right endpoint.

For \(b=\allowbreak 4\),
\(A=\allowbreak d+\xi,\allowbreak B=\allowbreak \xi,\allowbreak D_0=\allowbreak \delta_2,\allowbreak a_1=\allowbreak d^2+\delta_2\), and
\(3d^3\le\allowbreak  V\le\allowbreak  d(1/2+d)^2\). These give
\[
|\varepsilon|\le d/6,\quad\delta\le25d/24
\qquad(b=4,\ d\le1/4).
\tag{Q86.2.12}\label{eq:Q86:2.12}
\]
Indeed \(\varepsilon=\allowbreak -kd^2+k^2V/2\).
For its upper bound, the positive difference multiplied by \(m^2\) is
\((1-4d)(1+8d)/24\ge\allowbreak 0\).
For the lower bound let \(z=\allowbreak d/m\le\allowbreak 1/3\), so
\(\varepsilon/d\ge\allowbreak -z+3z^2/2\ge\allowbreak -1/6\).
The upper bound
\(\delta/d\le\allowbreak \theta/m+(1/2+d)^2/(2m^2)\) increases to \(25/24\).

For \(b=\allowbreak 3\), the parameters are
\[
\begin{gathered}
D_0=D_2=D_4=0,\quad D_1=-\delta_2,\quad D_3=\delta_3,\\
\quad a=k\delta_2,\quad\delta=k^2\delta_3/2,\\
\quad\eta=a-\delta\ge0,\quad\varepsilon=-\eta.
\end{gathered}
\tag{Q86.2.13}\label{eq:Q86:2.13}
\]
Since \(k\le\allowbreak 6/5,\allowbreak \delta_3\le\allowbreak \theta\delta_2\),
\(\delta/a\le\allowbreak 39/160,\allowbreak \eta\le\allowbreak 13/80<\allowbreak 1/6,\allowbreak \delta\le\allowbreak \eta/3,\allowbreak \delta\le\allowbreak 3d/25\).
The last uses \(\delta/d\le\allowbreak (18/25)\theta^2<\allowbreak 3/25\).

\paragraph{Q86.3 Finite propagation, actual losses and weighted AM--GM}
For every \(L\ge\allowbreak 2\), positivity, \(T(R^2)\le\allowbreak 2V\) and finite Duhamel expansion give
\[
\frac{T^LF_0}{m^2}\le
\bar f_L+kB_0^LR+\delta \bar f_{L-1}-a \bar f_{L-2}
+k\sum_{j=2}^{L-1}D_j\bar f_{L-1-j}.
\tag{Q86.3.1}\label{eq:Q86:3.1}
\]
The exact formula for any original \(R\) is
\(T^LR=\allowbreak B_0^LR+\sum_{j=\allowbreak 0}^{L-1}D_j\bar f_{L-1-j}\).
Insert \(F_0=\allowbreak m^2+mR+R^2/4\) and use
\(T^L(R^2)\le\allowbreak 2V\bar f_{L-1}\) to obtain \eqref{eq:Q86:3.1} term by term.
Define
\[
Z_L^*=S'_L+\delta S'_{L-1}-aS'_{L-2}+kB_0^LR
=\nu+Q_L(B_0)g',\qquad \alpha_L=\nu-EZ_L^*.
\]
Expanding by \eqref{eq:Q86:1.2}, the difference between \(Z_L^*\) and the right side of \eqref{eq:Q86:3.1} is
\(\sum_{j=\allowbreak 3}^L\gamma_j\bar f_{L-j}\), where
\[
\gamma_3=d-kD_2,\quad\gamma_4=d+\delta d-kD_3,
\quad
\gamma_j=\beta_j+\delta\beta_{j-1}-a\beta_{j-2}-kD_{j-1} (j\ge5).
\tag{Q86.3.2}\label{eq:Q86:3.2}
\]
Only existing nonnegative indices occur.
For \(L\ge\allowbreak 4\), \(\alpha_L=\allowbreak \gamma_{L+1}\).
The estimates used at every length are
\[
\begin{array}{c|c|c}
\text{range}&\gamma_j\ge hd\ (j\ge3)&\|Q_n\|\ (n\ge3)\\ \hline
b=3,\ d\le1/3&h=9/16-d/2&\le1\\
b=4,\ d\le1/4&h=1/2-6d/5&\le1+d\\
b\ge5,\ d\le1/5&h=1/2-6d/5&\le1+7d/5
\end{array}
\]
Here are the finite certificates for all indices.

For \(b=\allowbreak 3\),
\[
Q_n(z)=G_n(z)-\eta G_{n-2}(z)+\delta z^{n-2}-kz^{n+2}
\quad(n\ge3).
\tag{Q86.3.3}\label{eq:Q86:3.3}
\]
At \(z=\allowbreak i\sqrt\xi\), the real and normalized imaginary parts for \(n=\allowbreak 3\) are
\(1-\xi-\eta,\allowbreak 1+\delta-k\xi^2\).
Compared with \(P_3=\allowbreak G_3-kz^5\), the increase in squared modulus is at most
\(\eta[-2+(8/3)\xi+\eta(1+\xi/9)]\le\allowbreak -53\eta/108\),
where \(|P_3|\le\allowbreak 1\).
For \(n=\allowbreak 4\), the two components are nonnegative and bounded by those of
\(P_4=\allowbreak G_4-kz^6\); their lower bounds are \(11/36,\allowbreak 1/3\).
Moreover,
\(|P_4|\le\allowbreak (1-\xi)(1+\xi/2)+k\xi^3\le\allowbreak 1\), because
\(k\xi\le\allowbreak 39/80<\allowbreak 1/2\).

For every \(n\ge\allowbreak 5\),
\[
Q_n(z)=\frac{1-\eta+z^{n-2}[a-\delta z-z^2-kz^4+kz^5]}{1-z}.
\]
The last three terms have modulus at most \(\xi\), since
\((1-k\xi)^2+k^2\xi^3\le\allowbreak 1\), which follows from \(k\xi(1+\xi)\le\allowbreak 2\).
Also \(a+\delta\sqrt\xi\le\allowbreak 19\eta/12\) and
\((19/12)\xi^{3/2}<\allowbreak 1\), so the numerator has modulus at most
\(1+\xi^{5/2}\).
Finally \(2\xi^{3/2}+\xi^4\le\allowbreak 13/16<\allowbreak 1\), proving \(|Q_n|\le\allowbreak 1\).

For the losses, \(\gamma_3=\allowbreak d\) and
\(\gamma_4\ge\allowbreak (1-k\theta^2)d\ge\allowbreak 4d/5\).
For \(j\ge\allowbreak 5\), a positive \(D_{j-1}\) first occurs at \(D_7=\allowbreak \delta_5\), so
\(\gamma_j\ge\allowbreak [\kappa-k\theta d-k\theta^4]d\).
Subtracting \(h=\allowbreak 9/16-d/2\) and multiplying by \(m\) gives
\(1/32-\theta^4+(5/64)d-d^2/4\).
This is concave, and its two endpoint values are
\(4207/1048576\) and \(4207/1048576-1/576\), both positive.

For \(b\ge\allowbreak 4\), on the \(B_0\) frequency blocks write
\(R=\allowbreak (A-Bz)g'\), and
\[
P_n=G_n+kz^n(A-Bz),\qquad
Q_n=P_n+\varepsilon G_{n-2}+\delta z^{n-2}.
\tag{Q86.3.4}\label{eq:Q86:3.4}
\]
For \(b=\allowbreak 4\), use \(u_0=\allowbreak 1/2+d\); for \(b\ge\allowbreak 5\), use
\eqref{eq:Q86:2.7}. In both cases \(ku_0\le\allowbreak 1\).
For every \(n\ge\allowbreak 3\), \(\|P_n\|\le\allowbreak 1\):
at \(n=\allowbreak 3\), the components
\(1-\xi-k\xi^2B,\allowbreak 1-k\xi A\) are nonnegative and no larger than those of \(G_3\).
For \(n=\allowbreak 4\), the modulus is at most
\((1-\xi)(1+\xi/2)+\xi^2\le\allowbreak 1\).
For \(n=\allowbreak 5\), use
\(|G_5|\le\allowbreak 1-\xi/2+\xi^2/2\), then add
\(\xi^{5/2}\le\allowbreak 13\xi^2/20\).
The squared comparison has difference
\(\xi^2(1+2\xi-3\xi^2)/4\ge\allowbreak 0\).
For \(n\ge\allowbreak 6\), the modulus is at most
\((1+\xi^3)/\sqrt{1+\xi}+\xi^3\le\allowbreak 1\).
The last inequality is equivalent to
\(\xi(1+\sqrt{1+\xi})\le\allowbreak 1\), and this quantity is at most \(143/160\).

For \(Q_3\), use \(\|G_1\|=\allowbreak 1,\allowbreak \sqrt\theta<\allowbreak 13/20\).
The added cost for \(b=4\) is at most
\[
[1/6+(13/20)(25/24)]d=27d/32<d;
\]
for \(b\ge5\), it is at most
\[
[3/5+(13/20)(6/5)]d=69d/50<7d/5.
\]
For \(Q_4\), also use \(\|G_2\|\le\allowbreak 6/5\).
For larger \(n\), the geometric-sum norm is at most one and the power tail is smaller.
This proves the stated \(Q_n\) bounds.

For the losses, \(\gamma_3\ge\allowbreak  d,\allowbreak \gamma_4\ge\allowbreak \kappa d\).
When \(b=\allowbreak 4\), the same spectral estimate gives
\(\gamma_j\ge\allowbreak [\kappa-ku_0d-k\theta^2z_0]d\), \(j\ge\allowbreak 5\), with
\(z_0=\allowbreak \theta+11d/16\).
Subtract \(h\) and multiply by \(m\).
The resulting concave quadratic is
\(875/32768+(40369/81920)d-(11/5)d^2\),
whose endpoint values are \(875/32768,\allowbreak 4063/327680>\allowbreak 0\).
When \(b\ge\allowbreak 5\), every positive \(D_j\), \(j\ge\allowbreak 4\), is at most
\(\theta^2z_0d\), and \(a\le\allowbreak  d\).
Thus \(\gamma_j\ge\allowbreak [\kappa-d-k\theta^2z_0]d\).
After subtracting \(h\) and multiplying by \(m\), the difference is
\(875/32768-(3241/81920)d-9d^2/40\);
it decreases to \(40101/4096000>\allowbreak 0\).

For each \(L\ge\allowbreak 4\), retain the \(j=\allowbreak L,\allowbreak L-1\) losses in
\eqref{eq:Q86:3.2}. This gives
\[
0\le X_{b,L}/m^2\le Z:=Z_L^*-hd(1+q).
\tag{Q86.3.5}\label{eq:Q86:3.5}
\]
The finite polynomials satisfy \(Q_L=\allowbreak \nu+zQ_{L-1}\), hence
\[
Z=Bq-hd+N_0,\quad B=\nu-hd,\quad
N_0=B_0Q_{L-1}(B_0)g',\quad EN_0=-\alpha_L\le-hd.
\tag{Q86.3.6}\label{eq:Q86:3.6}
\]
In particular, the actual constant coordinate \(\nu-hd\) is retained.
The estimates below prove \(B>\allowbreak 0\) on each domain.
Weighted AM--GM between \(p\) and the nonnegative upper bound in
\eqref{eq:Q86:3.5} gives
\[
H_{b,L}\le m^2pZ
\le\left(c_*+N/2\right)^2,
\quad
c_*=m\sqrt B-\frac{mhd}{2\sqrt B},\quad
N=\frac m{\sqrt B}N_0.
\tag{Q86.3.7}\label{eq:Q86:3.7}
\]
The base is exactly \(m(\sqrt B\,\allowbreak p+Z/\sqrt B)/2\ge\allowbreak 0\).
If \(c_*\le\allowbreak  c\), then \(c+N/2\) remains nonnegative.
All row, second-moment and mean bounds for \(N\) come from this same original function.

For \(b=\allowbreak 3,\allowbreak L=\allowbreak 3\), the only retained loss is \(\gamma_3=\allowbreak d\).
The exact short-length formula is
\[
Z=(1-\eta)q-d+N_0,\quad N_0=B_0Q_2(B_0)g',
\quad Q_2=1+\delta+z-kz^4.
\tag{Q86.3.8}\label{eq:Q86:3.8}
\]
Here \(B=\allowbreak 1-\eta\), and \(hd\) in \eqref{eq:Q86:3.7} is replaced by \(d\);
\(\E N_0=\allowbreak -\alpha_3\).
The polynomial \(Q_2\) is estimated separately below.

\paragraph{Q86.4 Fixed-node Hermite--Lucas payment}
Let \(A,\allowbreak M>\allowbreak 0,\allowbreak \beta\ge\allowbreak 0\), and let the same real \(N\) satisfy
\[
|N|^2\le Md,\quad EN^2\le\theta Md,\quad EN\le-\beta d,
\quad 0\le H\le(c+N/2)^2,
\quad c=1-Ad>0,
\tag{Q86.4.1}\label{eq:Q86:4.1}
\]
Assume also \(c+N/2\ge\allowbreak 0\). Put
\(u=\allowbreak \sqrt{Md},\allowbreak x=\allowbreak u/(2c)<\allowbreak 1,\allowbreak k_0=\allowbreak 4A/M\) and \(\ell(d)=\allowbreak 2c\beta/M\).
For \(f(z)=\allowbreak (1+z/(2c))^{16}\), use quadratic Hermite interpolation with double contact at \(-\theta u\) and single contact at \(u\).
Since \(f'''\ge\allowbreak 0\) on \([-u,\allowbreak u]\), its remainder has sign
\((z+\theta u)^2(z-u)\le\allowbreak 0\).
Explicitly the majorant is
\(f(-\theta u)+f'(-\theta u)(z+\theta u)+B_H(z+\theta u)^2\), where
\[
B_H=
\frac{f(u)-f(-\theta u)-f'(-\theta u)(1+\theta)u}{[(1+\theta)u]^2}
\ge0.
\]
Its linear coefficient \(f'(-\theta u)+2\theta uB_H\) is also nonnegative.
Substituting the actual mean and second-moment bounds yields
\[
E(1+N/(2c))^{16}\le M_\theta(x)-\ell(d)xD_\theta(x),
\tag{Q86.4.2}\label{eq:Q86:4.2}
\]
\[
M_h(x)=\frac{h(1+x)^{16}+(1-hx)^{16}}{1+h},
\]
\[
D_h(x)=\frac{1-h}{1+h}16x(1-hx)^{15}
 +\frac{2h}{(1+h)^2}\{(1+x)^{16}-(1-hx)^{16}\}.
\]
Write the degree-\(j\) coefficients as
\(\binom{16}{j}M_j,\allowbreak \binom{16}{j}D_j\), respectively. Then
\[
M_j=\frac{h+(-h)^j}{1+h},\qquad
D_j=\frac{1-h}{1+h}j(-h)^{j-1}
 +\frac{2h}{(1+h)^2}(1-(-h)^j).
\tag{Q86.4.3}\label{eq:Q86:4.3}
\]
These \(M_j,\allowbreak D_j\) are scalar coefficients local to this paragraph, distinct from the original moments \(D_j\) in the propagation formulas.
Here \(D_0=\allowbreak 0\). At \(h=\allowbreak \theta\), the required low-degree data are
\[
\begin{gathered}
M_0=1,\quad M_1=0,\quad M_2=\theta,\qquad\\
560M_3=8645/64<136,\quad M_j\le M_4=10101/32768<5/16\ (j\ge4),
\end{gathered}
\]
\[
D_1=1,\quad D_2=0,\quad D_3=663/1024>16/25,
\quad560D_3>350,\quad D_j\ge0\ (j\ge0).
\tag{Q86.4.4}\label{eq:Q86:4.4}
\]
To see \(D_j\ge\allowbreak 0\), both terms are nonnegative when \(j\) is odd.
For even \(j\ge\allowbreak 4\), factor out \(h(1-h)/(1+h)^2\).
The remainder is
\(2\sum_{i=\allowbreak 0}^{j-1}h^i-jh^{j-2}(1+h)>\allowbreak 0\),
since \(0<\allowbreak h\le\allowbreak 1/2\) and \(jh^{j-2}\le\allowbreak 1\).
The cases \(j=\allowbreak 0,\allowbreak 1,\allowbreak 2\) are direct.

Let \(w=\allowbreak 1/c,\allowbreak v=\allowbreak w-1\ge\allowbreak 0\), so \(wv=\allowbreak k_0x^2\).
The finite Lucas identity is
\[
w^{16}+v^{16}=1+16t+104t^2+352t^3+660t^4
 +672t^5+336t^6+64t^7+2t^8,\quad t=wv.
\]
It follows from \(w-v=\allowbreak 1\) and the recursion
\(S_0=\allowbreak 2,\allowbreak S_1=\allowbreak 1,\allowbreak S_n=\allowbreak S_{n-1}+tS_{n-2}\) for
\(S_n=\allowbreak w^n+(-v)^n\).
Since \(v^{16}\le\allowbreak  t^8\),
\[
c^{-16}\ge1+16k_0x^2+104k_0^2x^4+352k_0^3x^6+660k_0^4x^8
 +672k_0^5x^{10}+336k_0^6x^{12}+64k_0^7x^{14}+k_0^8x^{16}.
\tag{Q86.4.5}\label{eq:Q86:4.5}
\]
If these coefficient bounds show
\[
c^{-16}-M_\theta(x)+\ell(d)xD_\theta(x)\ge r_*x^2,
\tag{Q86.4.6}\label{eq:Q86:4.6}
\]
then \eqref{eq:Q86:4.2} gives
\[
\boxed{EH^8\le1-(r_*M/4)c^{14}d.}
\tag{Q86.4.7}\label{eq:Q86:4.7}
\]

The three finite polynomials needed below are as follows.
Except for the first \(x^2,\allowbreak x^4\) coefficients, every even payment uses
\eqref{eq:Q86:4.5}.
For the \(x^3\) term, use \(560M_3<\allowbreak 136\);
for other odd terms, use \(M_j\le\allowbreak 5/16\), discarding positive \(D_j\) terms.
When \(k_0\ge\allowbreak 3\), after division of the difference by \(x^2\), the
\(x^4,\allowbreak x^6,\allowbreak x^8,\allowbreak x^{10},\allowbreak x^{12},\allowbreak x^{14}\) coefficients are at least
\(7000,\allowbreak 49000,\allowbreak 160000,\allowbreak 244000,\allowbreak 139000,\allowbreak 6500\).
The first is
\(352\cdot3^3-8008\cdot5/16=\allowbreak 14003/2>\allowbreak 7000\).
The others follow from
\(660\cdot3^4-12870\cdot5/16\),
\(672\cdot3^5-8008\cdot5/16\),
\(336\cdot3^6-1820\cdot5/16\),
\(64\cdot3^7-120\cdot5/16\), and \(3^8-5/16\).

If the first two coefficients are at least \(45/4,\allowbreak 629\), the difference divided by \(x^2\) is at least
\[
\begin{split}
P_A(x)={}&45/4-136x+629x^2-1365x^3+7000x^4-3575x^5
 +49000x^6-3575x^7\\
 &+160000x^8-1365x^9+244000x^{10}-175x^{11}
 +139000x^{12}-5x^{13}+6500x^{14}.
\end{split}
\]
Its full positive-quadratic decomposition is
\[
\begin{split}
P_A={}&(45/4-136x+420x^2)+x^2(209-1365x+2400x^2)\\
&+x^4(4600-3575x+1000x^2)+x^6(48000-3575x+100x^2)\\
&+x^8(159900-1365x+10x^2)+x^{10}(243990-175x+x^2)\\
&+x^{12}(138999-5x+x^2)+6499x^{14}.
\end{split}
\tag{Q86.4.8}\label{eq:Q86:4.8}
\]
Each quadratic has positive leading coefficient and negative discriminant.
For example, \(1365^2<\allowbreak 4\cdot209\cdot2400\) and
\(3575^2<\allowbreak 4\cdot4600\cdot1000\); the last four comparisons follow directly in the same way.
The first bracket has minimum \(101/420\), so \(P_A\ge\allowbreak 101/420\).

If \(k_0\ge\allowbreak 17/6\), the conservative high-even coefficient bounds are instead
\(5500,\allowbreak 38000,\allowbreak 120000,\allowbreak 170000,\allowbreak 90000,\allowbreak 4000\).
The first exceeds \(5500\) by \(209/54\); the others follow from
\[
\frac{4593655}{108}-\frac{32175}{8}>38000,\quad
\frac{9938999}{81}-\frac{5005}{2}>120000,
\]
\[
\frac{168962983}{972}-\frac{2275}{4}>170000,\quad
\frac{410338673}{4374}-\frac{75}{2}>90000,\quad
\frac{6975757441}{1679616}-\frac5{16}>4000.
\]
If the first two coefficients are at least \(15,\allowbreak 500\), use
\[
\begin{split}
P_B={}&15-136x+500x^2-1365x^3+5500x^4-3575x^5
 +38000x^6-3575x^7\\
&+120000x^8-1365x^9+170000x^{10}-175x^{11}
 +90000x^{12}-5x^{13}+4000x^{14}\\
={}&(15-136x+310x^2)+x^2(190-1365x+3000x^2)\\
&+x^4(2500-3575x+2000x^2)+x^6(36000-3575x+100x^2)\\
&+x^8(119900-1365x+10x^2)+x^{10}(169990-175x+x^2)\\
&+x^{12}(89999-5x+x^2)+3999x^{14}.
\end{split}
\tag{Q86.4.9}\label{eq:Q86:4.9}
\]
Every bracket is positive; the first has minimum \(13/155\).
Thus \(P_B\ge\allowbreak 13/155\).

Finally, if \(k_0\ge\allowbreak 3\) and the first two coefficients are at least \(35/4,\allowbreak 685\), use
\[
\begin{split}
P_C={}&(35/4-136x+530x^2)+x^2(155-1365x+3100x^2)\\
&+x^4(3900-3575x+1000x^2)+x^6(48000-3575x+100x^2)\\
&+x^8(159900-1365x+10x^2)+x^{10}(243990-175x+x^2)\\
&+x^{12}(138999-5x+x^2)+6499x^{14}.
\end{split}
\tag{Q86.4.10}\label{eq:Q86:4.10}
\]
Its expansion changes only the constant and \(x^2\) coefficients of \(P_A\), to \(35/4,\allowbreak 685\).
The second and third discriminants satisfy
\(1365^2<\allowbreak 4\cdot155\cdot3100\) and
\(3575^2<\allowbreak 4\cdot3900\cdot1000\); the remaining brackets are as in
\eqref{eq:Q86:4.8}.
The first has minimum \(27/1060\), so \(P_C\ge\allowbreak 27/1060\).
These decompositions hold for all real \(x\).

\paragraph{Q86.5 \(b=\allowbreak 3\) and every \(L\ge\allowbreak 4\)}
Use \(m=\allowbreak 1-d/2,\allowbreak \eta=\allowbreak k\delta_2-k^2\delta_3/2,\allowbreak h=\allowbreak 9/16-d/2\), and
\(B=\allowbreak 1-\eta-hd\).
From \(d\delta_3\ge\allowbreak \delta_2^2\), set
\(u=\allowbreak \delta_2/d,\allowbreak z=\allowbreak ku\le\allowbreak 39/80\). Then
\[
\eta/d\le z-z^2/2\le4719/12800<3/8.
\]
Also \(\delta/a\le\allowbreak 39/160\) gives
\[
\eta\ge(121/160)\delta_2\ge(121/160)d^2.
\tag{Q86.5.1}\label{eq:Q86:5.1}
\]
Therefore \(\eta+hd\le\allowbreak 15d/16-d^2/2\le\allowbreak 1-m^2\), so \(m^2\le\allowbreak  B\le\allowbreak 1\).
By \eqref{eq:Q86:3.6} and \(\|Q_{L-1}\|\le\allowbreak 1\), the same \(N\) in
\eqref{eq:Q86:3.7} satisfies
\[
|N|^2\le d,\qquad EN^2\le\theta d,\qquad
EN\le-\frac56\frac{19}{48}d=-\frac{95}{288}d.
\tag{Q86.5.2}\label{eq:Q86:5.2}
\]
The mean estimate uses \(m/\sqrt B\ge\allowbreak  m\ge\allowbreak 5/6\) and \(h\ge\allowbreak 19/48\).
Since \(B-hd/2>\allowbreak 0\) and \(m/\sqrt B\le\allowbreak 1\), use the whole expression
\[
c_*=(m/\sqrt B)(B-hd/2)
\le B-hd/2=1-\eta-3hd/2
\le1-27d/32-d^2/160\le c:=1-27d/32.
\tag{Q86.5.3}\label{eq:Q86:5.3}
\]
In the scalar lemma take \(M=\allowbreak 1,\allowbreak A=\allowbreak 27/32,\allowbreak \beta=\allowbreak 95/288\).
Then \(c\ge\allowbreak 23/32>\allowbreak 2/3,\allowbreak k_0=\allowbreak 27/8\), and
\(\ell\ge\allowbreak 2185/4608>\allowbreak 7/15\).
The first coefficients satisfy
\[
16(27/8)-120\theta+16(7/15)>45/4,
\]
\[
104(27/8)^2-1820(5/16)+350(7/15)>629.
\]
Since \(k_0\ge\allowbreak 3\), \(P_A\) pays all remaining terms.
By \eqref{eq:Q86:4.7} and \((2/3)^{14}>\allowbreak 1/300\),
\[
E H_{3,L}^8\le1-(101/1680)c^{14}d\le1-d/6000
\qquad(L\ge4).
\tag{Q86.5.4}\label{eq:Q86:5.4}
\]

\paragraph{Q86.6 The short boundary \(b=\allowbreak 3,\allowbreak L=\allowbreak 3\)}
Use \eqref{eq:Q86:3.8}, with
\(B=\allowbreak 1-\eta\ge\allowbreak 1-3d/8\), so again \(m^2\le\allowbreak  B\le\allowbreak 1\).
More precisely,
\[
m^2/B\le1-d/2,
\tag{Q86.6.1}\label{eq:Q86:6.1}
\]
because \((1-d/2)(1-3d/8)-m^2=\allowbreak d/8-d^2/16\ge\allowbreak 0\).
The exact \(Q_2=\allowbreak 1+\delta+z-kz^4\) needs its own estimate.
For \(\xi\in[0,\allowbreak \theta]\), \(1-k\xi^2\ge\allowbreak 0\), and
\[
|1+z-kz^4|^2=(1-k\xi^2)^2+\xi
\le1+\xi-2\xi^2+\xi^4
\le1+\xi-(11/6)\xi^2\le1+3/22<8/7.
\]
Here \(\theta^2<\allowbreak 1/6\).
Since \(\delta\le\allowbreak 3d/25\),
\(2\delta+\delta^2\le\allowbreak 153d/625<\allowbreak d/4\), whence
\[
|Q_2|^2\le8/7+d/4,\qquad
\frac{m^2}{B}\|Q_2\|^2
\le(1-d/2)(8/7+d/4)\le8/7<6/5.
\tag{Q86.6.2}\label{eq:Q86:6.2}
\]
On the other hand,
\(\alpha_3=\allowbreak d+\delta d-k\delta_3\ge\allowbreak 4d/5\).
Thus the same \(N=\allowbreak (m/\sqrt B)B_0Q_2g'\) satisfies
\[
|N|^2\le(6/5)d,\quad EN^2\le\theta(6/5)d,
\quad EN\le-(5/6)(4/5)d=-2d/3.
\tag{Q86.6.3}\label{eq:Q86:6.3}
\]
Because \(B\le\allowbreak 1\) and \(\sqrt B-d/(2\sqrt B)\) increases,
\[
c_*\le m(1-d/2)=m^2\le1-7d/8=:c.
\tag{Q86.6.4}\label{eq:Q86:6.4}
\]
Take \(M=\allowbreak 6/5,\allowbreak A=\allowbreak 7/8,\allowbreak \beta=\allowbreak 2/3\).
Then \(c\ge\allowbreak 17/24>\allowbreak 2/3,\allowbreak k_0=\allowbreak 35/12>\allowbreak 17/6\), and
\[
\ell=(10/9)c\ge85/108>39/50,\qquad
\ell\ge10/9-(175/54)x^2,
\tag{Q86.6.5}\label{eq:Q86:6.5}
\]
The last inequality uses \(d=\allowbreak 4c^2x^2/M\le\allowbreak (10/3)x^2\).
Apply the second bound in \eqref{eq:Q86:6.5} to \(D_1\), its first bound to \(D_3\), and discard the other nonnegative \(D_j\).
The first two coefficients are bounded below by
\[
16(35/12)-120\theta+16(10/9)=565/36>15,
\]
\[
104(35/12)^2-1820(5/16)
 +560(16/25)(39/50)-16(175/54)>500.
\]
The polynomial \(P_B\) therefore pays all higher terms, giving
\[
E H_{3,3}^8\le1-(39/1550)c^{14}d\le1-d/20000.
\tag{Q86.6.6}\label{eq:Q86:6.6}
\]
Here again \(c^{14}>\allowbreak (2/3)^{14}>\allowbreak 1/300\).
This proves every \(b=\allowbreak 3,\allowbreak L\ge\allowbreak 3\) row on the full domain.

\paragraph{Q86.7 \(b=\allowbreak 4,\allowbreak L\ge\allowbreak 4\) on \(0<\allowbreak d\le\allowbreak 1/4\)}
Now \(m=\allowbreak 1-d,\allowbreak h=\allowbreak 1/2-6d/5,\allowbreak B=\allowbreak 1+\varepsilon-hd\), with
\(|\varepsilon|\le\allowbreak  d/6\). Hence
\[
1-2d/3+6d^2/5\le B\le1,
\qquad B-m^2\ge4d/3+d^2/5>0.
\tag{Q86.7.1}\label{eq:Q86:7.1}
\]
Use \(\|Q_{L-1}\|\le\allowbreak 1+d\) and
\(m^2(1+d)^2=\allowbreak (1-d^2)^2\). We have
\[
\frac{17}{16}(1-2d/3+6d^2/5)-(1-d^2)^2
\ge\frac1{16}-\frac{17}{24}d+\frac{257}{80}d^2>0.
\tag{Q86.7.2}\label{eq:Q86:7.2}
\]
The estimate uses \(d^4\le\allowbreak  d^2/16\).
The final quadratic has negative discriminant because
\(289/576<\allowbreak 257/320\).
Thus the actual row and second moment of \(N\) can use \(M=\allowbreak 17/16\).
Also \(m/\sqrt B\ge\allowbreak  m\ge\allowbreak 3/4,\allowbreak h\ge\allowbreak 1/5\), so
\(\E N\le\allowbreak -3d/20\).

For the center, \(m\sqrt B\le\allowbreak (m^2+B)/2\) and \(m/\sqrt B\ge\allowbreak 3/4\) give
\[
c_*\le1-d+d^2/2+\varepsilon/2-(7/8)hd
\le1-(65/48)d+(31/20)d^2
\le1-(29/30)d=:c.
\tag{Q86.7.3}\label{eq:Q86:7.3}
\]
Take \(A=\allowbreak 29/30,\allowbreak \beta=\allowbreak 3/20\).
Then \(c\ge\allowbreak 91/120>\allowbreak 3/4,\allowbreak k_0=\allowbreak 928/255>\allowbreak 18/5\), and
\(\ell\ge\allowbreak 91/425>\allowbreak 1/5\).
The first coefficients obey
\[
16(18/5)-120\theta+16/5=241/20>45/4,
\]
\[
104(18/5)^2-1820(5/16)+350/5>629.
\]
The polynomial \(P_A\) and \eqref{eq:Q86:4.7} imply
\[
E H_{4,L}^8\le1-\frac{1717}{26880}c^{14}d
\le1-d/2000\qquad(L\ge4,\ d\le1/4).
\tag{Q86.7.4}\label{eq:Q86:7.4}
\]
Use \((3/4)^{14}>\allowbreak 1/64\).
The interval includes \(d=\allowbreak 1/4\).

\paragraph{Q86.8 Every \(b\ge\allowbreak 5,\allowbreak L\ge\allowbreak 4\) on \(0<\allowbreak d\le\allowbreak 1/5\)}
Here \(m=\allowbreak 1-9d/8,\allowbreak h=\allowbreak 1/2-6d/5\ge\allowbreak 13/50,\allowbreak B=\allowbreak 1+\varepsilon-hd\).
Equation \eqref{eq:Q86:2.11} gives
\[
B_{\rm low}:=1-11d/10+6d^2/5\le B
\le1+d/10+6d^2/5\le267/250,
\]
\[
B-m^2\ge(23/20)d-(21/320)d^2>0.
\tag{Q86.8.1}\label{eq:Q86:8.1}
\]
All these \(B\) are positive.
Using \(\|Q_{L-1}\|\le\allowbreak 1+7d/5\), calculate
\[
m^2(1+7d/5)^2=(1+11d/40-63d^2/40)^2
\le1+11d/20-(4919/1600)d^2.
\]
The remaining cubic and quartic terms are
\(d^3[-693/800+(3969/1600)d]\le\allowbreak 0\).
Therefore
\[
\frac54 B_{\rm low}-m^2(1+7d/5)^2
\ge\frac14-\frac{77}{40}d+\frac{7319}{1600}d^2>0,
\tag{Q86.8.2}\label{eq:Q86:8.2}
\]
The final discriminant comparison is \(5929<\allowbreak 7319\), so take \(M=\allowbreak 5/4\).
For the mean, use the precise ratio
\[
\frac m{\sqrt B}\ge\frac{149}{200},
\tag{Q86.8.3}\label{eq:Q86:8.3}
\]
Indeed \(m\ge\allowbreak 31/40,\allowbreak B\le\allowbreak 267/250\), and
\(961/1600-(149/200)^2(267/250)=\allowbreak 78583/10000000>\allowbreak 0\).
Thus \(\E N\le\allowbreak -(149/200)(13/50)d\le\allowbreak -19d/100\).

Using \(\sqrt B\le\allowbreak 1+(B-1)/2\) and \eqref{eq:Q86:8.3}, the center satisfies
\[
\begin{split}
c_*&\le(1-9d/8)(1+d/20+3d^2/5)-(149/400)hd\\
&\le1-(1009/800)d+(3963/4000)d^2
\le1-(21/20)d=:c.
\end{split}
\tag{Q86.8.4}\label{eq:Q86:8.4}
\]
The last comparison is
\(1009/800-3963/20000=\allowbreak 10631/10000>\allowbreak 21/20\).
Take \(A=\allowbreak 21/20,\allowbreak \beta=\allowbreak 19/100\).
Then \(c\ge\allowbreak 79/100>\allowbreak 3/4,\allowbreak k_0=\allowbreak 84/25\), and
\(\ell\ge\allowbreak 3002/12500>\allowbreak 6/25\).
The first coefficients satisfy
\[
16(84/25)-120\theta+16(6/25)=177/20>35/4,
\]
\[
104(84/25)^2-1820(5/16)+350(6/25)
=1723421/2500>685.
\]
The polynomial \(P_C\) pays the remaining terms, giving
\[
E H_{b,L}^8\le1-(27/3392)c^{14}d\le1-d/10000
\quad(b\ge5,L\ge4,d\le1/5).
\tag{Q86.8.5}\label{eq:Q86:8.5}
\]
Again use \((3/4)^{14}>\allowbreak 1/64\).
The scalar condition \(x<\allowbreak 1\) holds in all four applications:
\(c>\allowbreak 2/3,\allowbreak Md\le\allowbreak 5/12\), so \(x^2<\allowbreak 15/64\).
The nonnegative square base follows from the actual AM--GM comparison and \(c\ge\allowbreak  c_*\).

\paragraph{Q86.9 Actual terminal means and bare-path costs}
Let \(t_b=\allowbreak \E(h_2f_b)\).
Adjunction gives
\(t_b=\allowbreak \E(pf_{b+1})=\allowbreak 2\E f_{b+1}-\E f_{b+2}\).
Integrating \eqref{eq:Q86:1.2} also gives
\[
t_b=1-\beta_{b+1}-\beta_{b+2}+\beta_{b+3}
-\sum_{j=3}^b\beta_jt_{b-j}.
\tag{Q86.9.1}\label{eq:Q86:9.1}
\]
All \(t_j\ge\allowbreak 0\). The exact initial values are
\[
t_0=1+d,\quad t_1=1,\quad t_2=1-d-\delta_2,\quad t_3=1-2d-d^2,
\]
\[
t_4=1-3d+\delta_2-d^2+\delta_3,\qquad t_5=1-4d+2(1+d)\delta_2,
\]
\[
t_6=1-5d+2d^2+d^3+(3+2d)\delta_2-\delta_3-\delta_4.
\tag{Q86.9.2}\label{eq:Q86:9.2}
\]
Two further finite expansions give
\[
t_7=1-6d+5d^2+2d^3+4\delta_2-\delta_2^2-2(1+d)\delta_3,
\tag{Q86.9.3}\label{eq:Q86:9.3}
\]
\[
t_8=1-7d+9d^2+2d^3+(5-4d-3d^2)\delta_2-\delta_2^2
 -(3+2d)\delta_3+\delta_4+\delta_5.
\tag{Q86.9.4}\label{eq:Q86:9.4}
\]
For \(b=\allowbreak 7\), the base term is
\(1-d+\delta_2-\delta_3\), and the loss is
\(d(t_4+t_3)+\tau(t_2+t_1+t_0)+\delta_3t_0\).
Its exact expansion is
\(5d-5d^2-2d^3-3\delta_2+\delta_2^2+\delta_3+2d\delta_3\).
For \(b=\allowbreak 8\), the base is
\(1-d+\delta_2-\delta_3+\delta_4+\delta_5\), and the loss is
\(d(t_5+t_4)+\tau(t_3+t_2+t_1+t_0)+\delta_3(t_1+t_0)\).
It expands to
\(6d-9d^2-2d^3+(4d+3d^2-4)\delta_2+\delta_2^2+2\delta_3+2d\delta_3\).
Subtracting gives \eqref{eq:Q86:9.3}--\eqref{eq:Q86:9.4}.

Now assume \(1/5\le\allowbreak  d\le\allowbreak 1/3\).
For \(b=\allowbreak 5\), use \(\delta_2\le\allowbreak \theta d\):
\[
t_5-(1-d)/2\le1/2-(43/16)d+(13/16)d^2\le-1/200.
\tag{Q86.9.5}\label{eq:Q86:9.5}
\]
The quadratic decreases on the whole interval, with left endpoint \(-1/200\).
For \(b=\allowbreak 6\), \eqref{eq:Q86:9.2} similarly gives
\[
t_6-(1-d)/2\le1/2-(105/32)d+(45/16)d^2+d^3
\le-143/4000.
\tag{Q86.9.6}\label{eq:Q86:9.6}
\]
Its derivative at the right endpoint is still at most \(-103/96\), so it decreases throughout.
For \(b=\allowbreak 7\), discard two nonpositive terms and use \(\delta_2\le\allowbreak \theta d\), obtaining the convex upper bound
\[
t_7-(1-d)/2\le1/2-(31/8)d+5d^2+2d^3.
\]
Its endpoints are \(-59/1000,\allowbreak -35/216\).
For \(b=\allowbreak 8\),
\(\delta_4+\delta_5\le\allowbreak (\theta+\theta^2)\delta_3\le\allowbreak (3+2d)\delta_3\), and
\(5-4d-3d^2>\allowbreak 0\). Thus
\[
t_8-(1-d)/2\le1/2-(143/32)d+(59/8)d^2+(25/32)d^3.
\]
This is also convex, with endpoint values \(-37/400,\allowbreak -61/432\).

For every \(b\ge\allowbreak 9\), all \(\beta\) indices corresponding to
\(t_0,\allowbreak \ldots,\allowbreak t_6\) in \eqref{eq:Q86:9.1} are at least three.
Retain all seven nonnegative terms.
The finite alternating sums bound the base by
\(1-\beta_{b+1}-\beta_{b+2}+\beta_{b+3}\le\allowbreak 1-\tau+\delta_3\), hence
\[
t_b\le1-\tau(1+t_0+\cdots+t_6)+\delta_3,
\quad
t_0+\cdots+t_6=7-14d+d^3+(5+4d)\delta_2-\delta_4.
\tag{Q86.9.7}\label{eq:Q86:9.7}
\]
Let \(P\) be the quantity subtracted from one in \eqref{eq:Q86:9.7}.
Use \(\delta_4\le\allowbreak \theta^2\delta_2,\allowbreak \delta_3\le\allowbreak \theta\delta_2\), retaining the global
\(u=\allowbreak \delta_2/d\in[d,\allowbreak \theta]\):
\[
P\ge d[(1-u)(A+Cud)-\theta u],
\quad A=8-14d+d^3,\quad C=5+4d-\theta^2.
\tag{Q86.9.8}\label{eq:Q86:9.8}
\]
The bracket decreases with \(u\), since its derivative is
\(-A+Cd(1-2u)-\theta<\allowbreak 0\), using \(A\ge\allowbreak 10/3,\allowbreak Cd\le\allowbreak 19/9\).
Thus the global endpoint \(u=\allowbreak \theta\) gives
\[
P-(1+d)/2\ge-1/2+(8\kappa-\theta^2-1/2)d
-\kappa(14-5\theta+\theta^3)d^2.
\tag{Q86.9.9}\label{eq:Q86:9.9}
\]
Only nonnegative \(d^3,\allowbreak d^4\) terms were discarded here.
The coefficients satisfy
\(8\kappa-\theta^2-1/2=\allowbreak 4183/1024>\allowbreak 4\) and
\(\kappa(14-5\theta+\theta^3)=\allowbreak 7493391/1048576<\allowbreak 29/4\).
Therefore \eqref{eq:Q86:9.9} is at least
\(-1/2+4d-29d^2/4\).
This is concave, with endpoint values \(1/100,\allowbreak 1/36\), so it is at least \(1/100\).
Consequently \(t_b\le\allowbreak (1-d)/2-1/100\) for every \(b\ge\allowbreak 9\).
Together with the \(b=\allowbreak 5,\allowbreak 6,\allowbreak 7,\allowbreak 8\) estimates, this proves \eqref{eq:Q86:0.4}.
Since the original kernel of \(T\) is at most two,
\[
\boxed{0\le TF_b\le2t_b\le1-d\quad(b\ge5,d\ge1/5).}
\tag{Q86.9.10}\label{eq:Q86:9.10}
\]

For the bare-path cost, reserve the notation \(\mathcal V_n\) for the finite sum
\[
\mathcal V_n=\sum_{j=2}^n(-K)^j1
=B_0G_{n-1}(B_0)g'.
\]
Retaining \(\beta_n\bar f_0\) in \eqref{eq:Q86:1.2} gives
\(0\le\allowbreak \bar f_n\le\allowbreak  q+\mathcal V_n-\beta_n\), and
\(\E\mathcal V_n=\allowbreak -\beta_{n+1}\).
AM--GM for \(p\) and this nonnegative upper bound yields
\[
p \bar f_n\le[1+(\mathcal V_n-\beta_n)/2]^2,
\qquad E[(p \bar f_n)^8]\le e^{-8\beta_n}Ee^{8\mathcal V_n}.
\tag{Q86.9.11}\label{eq:Q86:9.11}
\]
The square base is nonnegative, so use \(y\le\allowbreak  e^{y-1}\) for \(y\ge\allowbreak 0\).
For \(n\ge\allowbreak 4\), \(\|G_{n-1}g'\|^2\le\allowbreak  d\).
By \eqref{eq:Q86:1.4} and
\(\beta_n,\allowbreak \beta_{n+1}\ge\allowbreak \kappa d\),
\[
E[(p \bar f_n)^8]\le\exp[-8\beta_n-8\beta_{n+1}+17d]
\le e^{(15/2)d}.
\tag{Q86.9.12}\label{eq:Q86:9.12}
\]
For \(n=\allowbreak 3\),
\(\mathcal V_3=\allowbreak r-s,\allowbreak \E\mathcal V_3=\allowbreak -d,\allowbreak \beta_3=\allowbreak d\).
Its predecessor is \(g'+B_0g'\), with squared norm exactly
\(d+\delta_2\le\allowbreak (1+\theta)d\); skew adjointness cancels the cross term.
Therefore
\[
E[(p \bar f_3)^8]\le\exp[(17(1+\theta)-16)d]=e^{253d/32}.
\tag{Q86.9.13}\label{eq:Q86:9.13}
\]
Since \(15/2\le\allowbreak 253/32\), these prove \eqref{eq:Q86:0.3}.
The \(n=\allowbreak 3\) calculation uses the actual \(G_2\) bound separately.

\paragraph{Q86.10 High domains and completion}
For \(b\ge\allowbreak 5,\allowbreak d\ge\allowbreak 1/5,\allowbreak L\ge\allowbreak 4\), the actual mean bound
\eqref{eq:Q86:9.10} and positivity give
\[
0\le X_{b,L}\le(1-d)\bar f_{L-1}.
\]
Since \(L-1\ge\allowbreak 3\), \eqref{eq:Q86:0.3} yields
\[
E H_{b,L}^8\le(1-d)^8e^{253d/32}
\le e^{-3d/32}\le1-d/11.
\tag{Q86.10.1}\label{eq:Q86:10.1}
\]
The last comparison uses
\(e^{-3d/32}\le\allowbreak (1+3d/32)^{-1}\) and \(d\le\allowbreak 1/3\).
This joins \eqref{eq:Q86:8.5} at \(d=\allowbreak 1/5\).

For \(b=\allowbreak 4,\allowbreak d\ge\allowbreak 1/4\), return to the actual \(t_4\) in
\eqref{eq:Q86:9.2}:
\[
t_4\le1-(2487/1024)d-d^2\le1353/4096<1/3,
\quad TF_4\le2/3.
\]
For every \(L\ge\allowbreak 4\), it follows that
\[
\log E H_{4,L}^8\le8\log(2/3)+(253/32)d
\le-28/9+253/96=-137/288\le-d.
\tag{Q86.10.2}\label{eq:Q86:10.2}
\]
Here \(\log(2/3)\le\allowbreak -7/18\) follows from
\(\log(1-x)\le\allowbreak -x-x^2/2\) at \(x=\allowbreak 1/3\).
Thus \(\E H_{4,\allowbreak L}^8\le\allowbreak  e^{-d}\le\allowbreak 1-3d/4\), joining
\eqref{eq:Q86:7.4} at \(d=\allowbreak 1/4\).
Sections Q86.5--Q86.6 give every \(b=\allowbreak 3,\allowbreak L\ge\allowbreak 3\) row on the full domain.
This proves \eqref{eq:Q86:0.2}, including \(d=\allowbreak 0\) and every stated integer boundary.

For decorations, positivity gives
\(0\le\allowbreak  E_rpT^L(E_tF_b)\le\allowbreak  pT^LF_b\).
Replace \(K\) by \(-K\) to reverse the terminal, connector, receiving unit arm and all decorations together. Then
\[
E\left[\left\{\overline E_rqA^L(\overline E_t\,\overline{h_2f_b})\right\}^8\right]
\le1-d/20000.
\]
The overline denotes the full reversal of the corresponding original edges.
